\subsubsection{Connected Autonomous Vehicle (CAV) Model}\label{CAV}

For connected autonomous vehicles, we adopt the point-mass longitudinal dynamics described in \cite{Mirabilio.2023}. The motion of vehicle $i$ is governed by

\begin{equation}
\dot{x}_i = v_i,
\qquad
\dot{v}_i = u_i,
\end{equation}
In a car-following configuration, the spacing and relative velocity are defined as
\begin{equation}
s_i = x_{i-1} - x_i - L,
\qquad
\Delta v_i = v_i - v_{i-1}.
\end{equation}

CAVs implement an Adaptive Cruise Control (ACC) strategy based on a Constant Time Headway (CTH) policy. The desired spacing is given by
\begin{equation}
s_i^{\mathrm{ref}} = s_0 + T_C v_i,
\end{equation}

where $s_0$ denotes the standstill spacing and $T_C > 0$ is the selected time headway.

Defining the spacing error
\begin{equation}
e_{p,i} = s_i - s_i^{\mathrm{ref}},
\end{equation}

its dynamics follow from
\begin{equation}
\dot{s}_i = -\Delta v_i,
\qquad
\dot{s}_i^{\mathrm{ref}} = T_C \dot{v}_i = T_C u_i,
\end{equation}

which yields
\begin{equation}
\dot{e}_{p,i} = -\Delta v_i - T_C u_i.
\end{equation}

A linear ACC controller that regulates both spacing and velocity errors can therefore be written as
\begin{equation}
u_i =
k_s ( s_i - s_i^{\mathrm{ref}} )
-
k_v \Delta v_i,
\end{equation}

where $k_s > 0$ and $k_v > 0$ denote the spacing and velocity feedback gains, respectively. This control law constitutes a linear realization of the CTH policy and is consistent with the mesoscopic ACC formulation proposed in \cite{Mirabilio.2023}.

An important property of the CTH-based ACC strategy is its capability to guarantee weak string stability using only local measurements \cite{Swaroop.2001}. In the frequency domain, the disturbance propagation transfer function satisfies
\begin{equation}
\| H(j\omega) \|_{\infty} \le 1,
\end{equation}
ensuring that spacing disturbances are not amplified upstream along a vehicle string.
Moreover, robustness with respect to sensing and actuation delays imposes a constraint on the admissible implementation lag. In particular, the total delay must satisfy
\begin{equation}
T_C \le \frac{h_w}{2},
\end{equation}
indicating that allowable delay is directly limited by the chosen time headway \cite{Swaroop.2001}.